\chapter{Conclusion and Future Work}
\label{ch:conclusion}

The three research questions of \cref{ch:intro} can now be answered
directly.

\emph{Does learned attention add value over unlearned relational
propagation?} On the equity track, yes, with unusual consistency for a
result in this domain: positive attention value-add in thirty of thirty
seeded runs, across two hyperparameter configurations and two devices,
against an anchor that shares the model's inputs and topology. The
mechanism is modest and now understood --- a stationary, near-uniform
reweighting of neighbour pooling --- and the resulting composite is a
real, unique, consistent factor that nevertheless does not beat the best
single alpha it was built from. On the energy forecasting track the
answer is narrower: attention improves cross-sectional ranking over the
uniform anchor in five of five seeds but not mean-squared skill, where
uniform pooling is already near-optimal.

\emph{Does one kernel transfer across heterogeneous graphs without
manufacturing false positives?} Yes, and the negative controls are the
evidence: the same code that finds a consistent positive margin on real
equity data finds nothing on synthetic random walks, finds nothing on
synthetic uncoupled power zones, and --- when real energy data produced
numbers too good to be true --- the surrounding discipline treated the
magnitude as an alarm rather than a result. The transfer claim is
literal: one model class, one training loop, one gate suite, two
structurally different graphs.

\emph{Does the physical topology improve forecasts?} Yes, where it
should. The interconnector graph lifts node-level price forecast skill
by $+0.131$ over an identically-informed no-graph baseline, and the
clearest relational gain sits on the irreducibly relational target:
edge-level spread prediction beats a both-endpoint baseline by $+0.056$
skill in five of five seeds, while a synthetic control inverts the
comparison. Congestion as an explicit edge feature added nothing robust
under either of two operationalisations --- an honest null.

Beyond the questions it set out to answer, the project's most
transferable output is methodological: evaluate on untransformed
returns; treat rank correlation and profitability as different claims;
treat implausible magnitudes as alarms; require seeds, fixed devices,
and cross-implementation replication before believing a learned effect;
and keep negative controls running permanently, because they are what
give positive results their meaning.

\section*{Future work}

Four directions follow directly from the open ends documented in the
experiment log and forecasting notes, in decreasing order of expected
value. First, data: a survivorship-bias-free equity vendor would turn
the biased absolute Sharpe levels into defensible ones, and longer,
multi-window energy samples would test the six-month forecasting results
out-of-sample in time. Second, evaluation: the value-added gate deserves
its portfolio-level variant --- the marginal contribution of the
composite to a multifactor portfolio --- alongside the max-of-singles
bar it never cleared. Third, the forecasting track is untuned: a
hyperparameter grid, seed ensembles, and walk-forward retraining, all
already built for the alpha track, transfer mechanically; an edge-level
congestion target (not just the spread) is the natural next prediction
problem. Fourth, the congestion hypothesis remains open rather than
dead: the flow-based-coupling regions publish shadow prices on critical
network elements, and a congestion signal built from those --- rather
than from flow-over-capacity ratios with partial coverage --- is the
experiment this project would run next. Any renewed energy \emph{alpha}
attempt, by contrast, should wait for genuinely tradeable targets:
balancing or intraday spreads, or transmission-rights-aware profit
accounting.

The report closes where it began. The coupling structure of markets is
information; on the evidence here, learning how to read it earns a
consistent but modest premium over reading it naively, and the
discipline required to measure that premium honestly --- anchors,
controls, embargoes, seeds, and the willingness to let three spectacular
numbers die --- is itself the result this project most confidently
recommends to its successors.
